% !PRINTED 322
\weberpage{322}
%<<CONT
ish if (4) has two equal roots, and by XIV this is possible only
when $f(x,y)$ is the $n$th power of a linear function.
In this case the determinant (5) vanishes identically: if we
therefore set this case aside, it follows \gesperrt{that for an equation with only
real roots, the equation obtained by setting the Hesse
determinant equal to zero has only imaginary roots.}

If we wish to return to the inhomogeneous representation, we can, making
use of the Euler theorem
\weq{6}{\begin{aligned}
xf'(x) + yf'(y) &= nf(x)\\
xf''(x,x) + yf''(x,y) &= (n-1)f'(x)\\
xf''(y,x) + yf''(y,y) &= (n-1)f'(y)
\end{aligned}}
eliminate $f'(y), f''(x,y), f''(y,y)$ by means of it. One obtains in this way,
on setting $y = 1$,
\weq{7}{\begin{aligned}
f'(y) &= nf(x) - xf'(x)\\
f''(x,y) &= (n-1)f'(x) - xf''(x)\\
f''(y,y) &= n(n-1)f(x) - 2(n-1)xf'(x) + x^2 f''(x),
\end{aligned}}
and hence, if
\weq{8}{(n-1)H(x,y) = f''(x,x)f''(y,y) - f''(x,y)^2}
is set:
\weq{9}{H(x,y) = H(x) = nf(x)f''(x) - (n-1)f'(x)^2.}
$H(x)$ is, with respect to $x$, at most of degree $2n-4$.

\begin{center}
\S.\,107.

The theorems of Laguerre for equations with only real roots.
\end{center}

Of the theorems last derived, \gesperrt{Laguerre} has made an application
to equations with only real roots, which we shall now become acquainted
with.

Let $f(x) = 0$ be an equation of degree $n$ with $n$ \gesperrt{real}
roots $\alpha, \alpha_1, \alpha_2 \dots$, \gesperrt{which we assume to be distinct
from one another}.

If $x$ is a variable quantity, we have the formula, already applied
several times [\S.\,14, (11)]:
%CONT>>
